\subsection{Classical Language Models: n-grams}
\label{sec:ngram-models}
\slidesrc{Language models \& RNNs: 12--34}

Before neural networks, language models were built by \textbf{counting}. This section
shows how, why it works surprisingly well, and why it eventually hits a wall.

\subsubsection{n-grams and the Markov assumption}
\slidesrc{Language models \& RNNs: 13--16}

\begin{nbdef}{n-gram}
An \textbf{n-gram} is a sequence of $n$ consecutive words of a text.
\end{nbdef}

For the sentence ``The quick brown fox jumps over the lazy dog'':
\begin{center}
\small
\begin{tabular}{@{}ll@{}}
\toprule
$n$ & n-grams \\
\midrule
1 (\textbf{unigrams})  & \emph{the}, \emph{quick}, \emph{brown}, \emph{fox}, \ldots \\
2 (\textbf{bigrams})   & \emph{the quick}, \emph{quick brown}, \emph{brown fox}, \emph{fox jumps}, \ldots \\
3 (\textbf{trigrams})  & \emph{the quick brown}, \emph{quick brown fox}, \emph{brown fox jumps}, \ldots \\
4 (\textbf{four-grams})& \emph{the quick brown fox}, \emph{quick brown fox jumps}, \ldots \\
\bottomrule
\end{tabular}
\end{center}

An \textbf{n-gram language model}:
\begin{enumerate}[nosep]
  \item collects the \textbf{frequencies} (counts) of all n-grams in a large text
        corpus;
  \item uses these frequencies to predict the next word.
\end{enumerate}

\paragraph{The Markov assumption.}
The chain rule (Section~\ref{sec:what-is-a-language-model}) needs the probability
of the next word given the
\emph{whole} history. That is impossible to estimate by counting: most long histories
appear only once in any corpus, or never. An n-gram model makes a simplifying
assumption, named after the mathematician Markov (we met the same idea in Hidden Markov
Models, Section~\ref{sec:hmm-pos-tagging}).

\begin{nbdef}{Markov assumption (n-gram model)}
The next word depends only on the $n-1$ words just before it:
\[
  p\bigl(\vect x^{(t+1)} \mid \vect x^{(t)}, \ldots, \vect x^{(1)}\bigr) \approx
  p\bigl(\vect x^{(t+1)} \mid
  \underbrace{\vect x^{(t)}, \ldots, \vect x^{(t-n+2)}}_{n-1 \text{ words}}\bigr).
\]
(Counting the context: from position $t-n+2$ to position $t$ there are
$t - (t-n+2) + 1 = n-1$ words. Together with the predicted word they form one n-gram.)
\end{nbdef}

For example, a \textbf{bigram} model ($n = 2$) looks only at the previous word. The
probability of ``the quick brown fox'' is decomposed as
\[
  p(\texttt{<s>}) \, p(\textit{the} \mid \texttt{<s>}) \, p(\textit{quick} \mid \textit{the})
  \, p(\textit{brown} \mid \textit{quick}) \, p(\textit{fox} \mid \textit{brown}) \,
  p(\texttt{</s>} \mid \textit{fox}).
\]
\texttt{<s>} and \texttt{</s>} are two special tokens marking the \textbf{beginning} and
the \textbf{end of the sentence}. They let the model learn which words tend to
\emph{start} a sentence ($p(\textit{the} \mid \texttt{<s>})$ is high) and which tend to
\emph{end} it. Every sentence starts with \texttt{<s>}, so $p(\texttt{<s>}) = 1$.

\paragraph{Estimating the probabilities by counting.}
By the definition of conditional probability, the probability we need is a ratio:
\[
  p\bigl(\vect x^{(t+1)} \mid \vect x^{(t)}, \ldots, \vect x^{(t-n+2)}\bigr) =
  \frac{p\bigl(\vect x^{(t+1)}, \vect x^{(t)}, \ldots, \vect x^{(t-n+2)}\bigr)}
       {p\bigl(\vect x^{(t)}, \ldots, \vect x^{(t-n+2)}\bigr)} .
\]
We do not know these probabilities, but we have a corpus, a finite sample of language.
The \textbf{finite-sample approximation} replaces each probability with how often the
sequence occurs in the corpus:
\[
  p\bigl(\vect x^{(t+1)} \mid \vect x^{(t)}, \ldots, \vect x^{(t-n+2)}\bigr) \approx
  \frac{\overbrace{\text{count}\bigl(\vect x^{(t+1)}, \vect x^{(t)}, \ldots,
        \vect x^{(t-n+2)}\bigr)}^{\text{an n-gram}}}
       {\underbrace{\text{count}\bigl(\vect x^{(t)}, \ldots,
        \vect x^{(t-n+2)}\bigr)}_{\text{an $(n-1)$-gram}}} .
\]
In words: \emph{among all the times the context appeared, in what fraction was it
followed by this word?} The quality of the approximation depends on the
\textbf{size of the corpus}: the more text, the closer the fractions get to the true
probabilities.

\begin{nbexample}{a four-gram model}
\slidesrc{Language models \& RNNs: 17--19}
We want to continue ``Anna was tired so she went to \underline{\quad}'' with a
\textbf{four-gram} model ($n = 4$). It looks only at the last $n - 1 = 3$ words,
\emph{she went to}, and discards the rest:
\[
  p(x \mid \textit{she went to}) \approx
  \frac{\text{count}(\textit{she went to } x)}{\text{count}(\textit{she went to})} .
\]
Suppose the corpus contains:
\begin{center}
\small
\begin{tabular}{@{}lrl@{}}
\toprule
sequence & count & estimated probability of $x$ \\
\midrule
\emph{she went to}        & 1000 & (the context) \\
\emph{she went to the}    & 500  & $500/1000 = 0.50$ \\
\emph{she went to work}   & 250  & $250/1000 = 0.25$ \\
\emph{she went to school} & 200  & $200/1000 = 0.20$ \\
\emph{she went to bed}    & 50   & $50/1000 = 0.05$ \\
\bottomrule
\end{tabular}
\end{center}
The four probabilities add up to 1 (here every occurrence of the context is followed by
one of these four words). The model predicts \emph{the} as most likely and \emph{bed} as
least likely.

\textbf{Was it a good idea to discard the longer context?} No: ``Anna was
\emph{tired}'' is exactly the clue that makes \emph{bed} the most sensible
continuation, but the model never saw it, because it fell outside the window of three
words.
\end{nbexample}

\begin{nbexample}{a bigram model from a toy corpus}
\notinslides
The corpus has three sentences:
\begin{center}
\texttt{<s>} I like cats \texttt{</s>} \qquad
\texttt{<s>} I like dogs \texttt{</s>} \qquad
\texttt{<s>} you like cats \texttt{</s>}
\end{center}
Counts we need: \texttt{<s>} appears 3 times, followed twice by \emph{I} and once by
\emph{you}; \emph{I} appears 2 times, always followed by \emph{like}; \emph{like}
appears 3 times, followed twice by \emph{cats} and once by \emph{dogs}; \emph{dogs}
appears once, followed by \texttt{</s>}. So
\[
  p(\textit{I} \mid \texttt{<s>}) = \tfrac23, \quad
  p(\textit{like} \mid \textit{I}) = \tfrac22 = 1, \quad
  p(\textit{dogs} \mid \textit{like}) = \tfrac13, \quad
  p(\texttt{</s>} \mid \textit{dogs}) = \tfrac11 = 1 .
\]
The probability of the sentence ``I like dogs'' is
\begin{align*}
  p(\textit{I like dogs})
  &= p(\textit{I} \mid \texttt{<s>})\, p(\textit{like} \mid \textit{I})\,
     p(\textit{dogs} \mid \textit{like})\, p(\texttt{</s>} \mid \textit{dogs}) \\
  &= \tfrac23 \cdot 1 \cdot \tfrac13 \cdot 1 = \tfrac29 \approx 0.22 .
\end{align*}
Two observations that lead to the next topic:
\begin{itemize}[nosep]
  \item ``you like dogs'' \emph{never} appears in the corpus, yet the model gives it
        $\frac13 \cdot 1 \cdot \frac13 \cdot 1 = \frac19$: by cutting sentences into
        short pieces, the model \textbf{generalises} to new sentences;
  \item ``I hate cats'' gets probability \textbf{0}, because the bigram \emph{I hate}
        was never seen. One unseen bigram is enough to kill the whole sentence.
\end{itemize}
\end{nbexample}

\subsubsection{Sparsity problems}
\slidesrc{Language models \& RNNs: 20--26}

\paragraph{Sparsity problem 1: unseen continuations.}
For many words $x$, the sequence ``she went to $x$'' \textbf{never occurs} in the
corpus (``she went to Reykjavik'', ``she went to dinner'' may simply be absent). Their
count is 0, so their probability is 0: the model says they are \emph{impossible}. Many
words with probability 0 make the LM \textbf{inflexible}.

\textbf{Mitigation: smoothing.} Add a small value $\epsilon$ to the count of every
word $x \in V$, so that nothing is ever exactly 0:
\[
  p(x \mid \textit{she went to}) =
  \frac{\text{count}(\textit{she went to } x) + \epsilon}
       {\text{count}(\textit{she went to}) + \epsilon |V|} .
\]

\begin{nbwarning}[the denominator must change too]
The slide adds $\epsilon$ only to the numerator. Then the probabilities no longer add
up to 1: summing the numerators over all $|V|$ words gives
$\text{count}(\textit{she went to}) + \epsilon|V|$, not
$\text{count}(\textit{she went to})$. Adding $\epsilon|V|$ to the denominator, as
above, fixes this. With $\epsilon = 1$ the method is called \textbf{add-one} or
\textbf{Laplace smoothing}.
\end{nbwarning}

\begin{nbexample}{add-one smoothing on the toy corpus}
\notinslides
In the toy corpus above, a word can be followed by one of the 6 tokens
\emph{I}, \emph{you}, \emph{like}, \emph{cats}, \emph{dogs}, \texttt{</s>}, so
$|V| = 6$; and $\text{count}(\textit{I}) = 2$. With $\epsilon = 1$:
\[
  p(\textit{like} \mid \textit{I}) = \frac{2 + 1}{2 + 6} = \frac38, \qquad
  p(\textit{cats} \mid \textit{I}) = \frac{0 + 1}{2 + 6} = \frac18 ,
\]
and the same $\frac18$ for each of the other 4 unseen tokens.
Check: $\frac38 + 5 \cdot \frac18 = 1$. \emph{like} is still by far the most likely, but
no word is impossible any more. The price: probability mass has been taken from the
words that were actually seen ($1 \to \frac38$) and spread over unseen ones.
\end{nbexample}

\paragraph{Sparsity problem 2: unseen contexts.}
What if ``she went to'' \textbf{never appeared} in the corpus? Then the denominator is
0, and we cannot compute the probability at all.

\textbf{Mitigation: backoff.} Condition on a shorter context, just ``went to'' (or
``she went'', as the slide writes), and if that is also missing, shorter still, down to
the unigram probability of $x$. Shorter contexts are more frequent, so they are more
likely to have been observed.

\begin{nbfact}{Increasing n makes sparsity worse}
A longer n-gram is a more specific sequence, so it is seen more rarely. With a larger
$n$ the model uses more context, but more of its counts are 0 and smoothing/backoff
are needed more often. In practice $n \le 5$.
\end{nbfact}

\paragraph{An example of sparsity.}
The \textbf{Berkeley Restaurant Project} dataset contains 9332 sentences of people
asking about restaurants, with a vocabulary of $|V| = 1446$ words. The table shows the
bigram counts for 8 of these words: the entry in row \emph{want}, column \emph{to} is
how many times \emph{want} was immediately followed by \emph{to}.

\begin{center}
\small
\begin{tabular}{@{}l*{8}{r}@{}}
\toprule
 & \emph{i} & \emph{want} & \emph{to} & \emph{eat} & \emph{chinese} & \emph{food} &
   \emph{lunch} & \emph{spend} \\
\midrule
\emph{i}       & 5  & 827 & 0  & 9   & 0 & 0 & 0 & 2   \\
\emph{want}    & 2  & 0  & 608 & 1   & 6  & 6  & 5  & 1   \\
\emph{to}      & 2  & 0  & 4   & 686 & 2  & 0 & 6  & 211 \\
\emph{eat}     & 0 & 0  & 2   & 0  & 16 & 2  & 42 & 0  \\
\emph{chinese} & 1  & 0  & 0  & 0  & 0 & 82 & 1  & 0  \\
\emph{food}    & 15 & 0  & 15  & 0  & 1  & 4  & 0 & 0  \\
\emph{lunch}   & 2  & 0  & 0  & 0  & 0 & 1  & 0 & 0  \\
\emph{spend}   & 1  & 0  & 1   & 0  & 0 & 0 & 0 & 0  \\
\bottomrule
\end{tabular}
\captionof{table}{Bigram counts in the Berkeley Restaurant Project (row = first word,
column = second word). Half of the 64 entries are 0, even among frequent words.}
\label{tab:berkeley-bigrams}
\end{center}

32 of the 64 cells are zero. Some zeros are real impossibilities (\emph{want want}),
but others are perfectly good English that simply did not occur (\emph{lunch food}?
\emph{spend lunch}?). And this is with \textbf{bigrams}: there are
$1446^2 \approx 2.1$ million possible bigrams, far more than a corpus of 9332 sentences
can show. With $n > 2$ it is much worse.

\subsubsection{Storage problems}
\slidesrc{Language models \& RNNs: 27}

An n-gram model must \textbf{store the counts} of all n-grams seen in the corpus. As $n$
grows, or as the corpus grows, this can become huge. With $|V| = 1446$ there are
\[
  1446^4 \approx 4.37 \times 10^{12}
\]
possible four-grams: naively storing a count for each of them, even at just one byte per
count, requires about \textbf{4 TB} of memory, for a toy vocabulary. Real vocabularies
have tens of thousands of words. (Storing only the n-grams actually seen saves a lot,
but on a large corpus that is still billions of entries.)

\subsubsection{Generating text with an n-gram model}
\slidesrc{Language models \& RNNs: 28--34}

A language model can also \textbf{generate} text: start from a context, compute the
distribution of the next word, \textbf{sample} a word from it, append it to the context
and repeat.

\begin{nbdef}{Sampling from a distribution}
\notinslides
\textbf{Sampling} a word means drawing it at random \emph{with the probabilities given
by the model}: a word with probability 0.3 is picked, on average, 3 times out of 10. It
is like rolling a die whose faces are the words of the vocabulary and whose faces have
different sizes. Sampling (instead of always taking the most likely word) gives varied
text rather than the same sentence every time.
\end{nbdef}

The slides use a \textbf{trigram} model built on the \textbf{Reuters corpus} (1.7
million words of financial news). Starting from ``today the'':

\begin{center}
\begin{tikzpicture}
\begin{axis}[xbar, width=6.5cm, height=4.2cm, bar width=7pt, xmin=0, xmax=0.2,
    symbolic y coords={emirate,italian,price,bank,company}, ytick=data,
    y tick label style={font=\small\itshape}, x tick label style={font=\scriptsize,
    /pgf/number format/fixed}, nodes near coords, nodes near coords style={font=\scriptsize},
    every node near coord/.append style={/pgf/number format/fixed,
    /pgf/number format/precision=3},
    title={\small $p(x \mid \textit{today the})$}, axis x line*=bottom, axis y line*=left,
    enlarge y limits=0.15]
  \addplot[fill=nbblue!60, draw=nbblue] coordinates
    {(0.153,company) (0.153,bank) (0.077,price) (0.039,italian) (0.039,emirate)};
\end{axis}
\end{tikzpicture}
\captionof{figure}{The five most likely words after ``today the'' in the trigram model
(the rest of the probability is spread over many other words).}
\end{center}

\begin{enumerate}[nosep]
  \item Context \emph{today the} $\to$ distribution: \emph{company} 0.153,
        \emph{bank} 0.153, \emph{price} 0.077, \emph{italian} 0.039,
        \emph{emirate} 0.039, \ldots\ Sample: \emph{price}.
  \item Text \emph{today the price}; context (last 2 words) \emph{the price} $\to$
        \emph{of} 0.308, \emph{for} 0.050, \emph{it} 0.046, \emph{to} 0.046,
        \emph{is} 0.031, \ldots\ Sample: \emph{of}.
  \item Text \emph{today the price of}; context \emph{price of} $\to$ \emph{the} 0.072,
        \emph{18} 0.043, \emph{oil} 0.043, \emph{its} 0.036, \emph{gold} 0.018, \ldots
        and so on.
\end{enumerate}
Note that at each step only the last $n - 1 = 2$ words are used as context: the
beginning \emph{today} is forgotten as soon as it leaves the window.

\paragraph{The effect of $n$.}
The slide shows text generated by n-gram models trained on a novel (Mary Shelley's
\emph{Frankenstein}), for increasing $n$:
\begin{center}
\small
\begin{tabularx}{\linewidth}{@{}c X@{}}
\toprule
$n$ & generated text \\
\midrule
1 & ``care I of . its destined , in , from . . for the Felix an not which measure excited'' \\
2 & ``I compassionated him , how heavily ; but I had arrived , but they averred , my unhallowed damps and'' \\
3 & ``I continued their single offspring . At length I gathered from a man who , born in freedom , spurned'' \\
4 & ``I continued walking in this manner , during which I enjoyed the feeling of happiness . Still thou canst listen'' \\
5 & ``I continued walking in this manner for some time , and I feared the effects of the dæmon's disappointment .'' \\
\bottomrule
\end{tabularx}
\end{center}
\begin{itemize}[nosep]
  \item Increasing $n$, the model is \textbf{more coherent}: unigrams give a word salad,
        bigrams give locally plausible pairs, 4- and 5-grams read like real sentences.
  \item Increasing $n$ \textbf{increases the sparsity problem}.
  \item With a high $n$ the model \textbf{repeats the training set}: most 5-word
        contexts were seen only once, so they have only one possible continuation, and
        the model copies whole passages of the novel instead of creating new text.
\end{itemize}

\begin{nbsummary}[n-gram language models]
\begin{itemize}[nosep]
  \item Markov assumption: the next word depends only on the previous $n-1$ words.
  \item Estimate: $p(x \mid \text{context}) \approx
        \text{count}(\text{context}, x) / \text{count}(\text{context})$.
  \item Sparsity 1 (unseen n-gram, probability 0) $\to$ smoothing; sparsity 2 (unseen
        context, division by 0) $\to$ backoff. Both get worse as $n$ grows; $n \le 5$.
  \item Storage grows with $n$ and with the corpus.
  \item Small $n$: incoherent text; large $n$: coherent but copied from the corpus.
  \item The context is cut off: anything older than $n-1$ words is forgotten.
\end{itemize}
\end{nbsummary}

\begin{nbexercise}{1}{bigram probabilities}
\notinslides
Using the toy corpus of this section (\texttt{<s>} I like cats \texttt{</s>},
\texttt{<s>} I like dogs \texttt{</s>}, \texttt{<s>} you like cats \texttt{</s>}) and a
bigram model without smoothing, compute $p(\textit{you like cats})$. Which sentence of
the form ``\emph{subject} like \emph{animal}'' has the highest probability?
\end{nbexercise}
\begin{nbsolution}
Counts: $p(\textit{you} \mid \texttt{<s>}) = \frac13$,
$p(\textit{like} \mid \textit{you}) = \frac11 = 1$,
$p(\textit{cats} \mid \textit{like}) = \frac23$,
$p(\texttt{</s>} \mid \textit{cats}) = \frac22 = 1$. So
\[
  p(\textit{you like cats}) = \tfrac13 \cdot 1 \cdot \tfrac23 \cdot 1 = \tfrac29 .
\]
The four candidates are
$p(\textit{I like cats}) = \frac23 \cdot \frac23 = \frac49$,
$p(\textit{I like dogs}) = \frac23 \cdot \frac13 = \frac29$,
$p(\textit{you like cats}) = \frac13 \cdot \frac23 = \frac29$,
$p(\textit{you like dogs}) = \frac13 \cdot \frac13 = \frac19$ (the factors equal to 1
are omitted). Check: they add up to $\frac{4+2+2+1}{9} = 1$, since these are the only
sentences the model can produce. \emph{I like cats} wins: it combines the most
frequent start (\emph{I}) and the most frequent animal (\emph{cats}).
\end{nbsolution}
